\documentclass[11pt,a4paper]{article}

\usepackage[T1]{fontenc}
\usepackage[utf8]{inputenc}
\usepackage[spanish]{babel}
\usepackage[a4paper,margin=2.8cm]{geometry}
\usepackage{microtype}
\usepackage[hidelinks]{hyperref}

\hypersetup{
  pdftitle={Menelao}
}

\setlength{\parindent}{1.25em}
\setlength{\parskip}{0.35em}
\raggedbottom

\title{Menelao}
\author{}
\date{}

\begin{document}

\maketitle

\begin{verse}
I

Uno de los mejores momentos del año\\
es cuando viajas a Esparta.\\
Menelao organiza una reunión, a la que acuden\\
los que sobrevivieron a Troya y al Nostos.

Por supuesto, eres el invitado de honor,\\
y te agasajan con comida y vino,\\
y hay una linda esclava de manos suaves\\
que cuida bien de tu cuerpo cansado.

Y todos lo pasáis bien, recordando\\
episodios repetidos tantas veces con los años\\
que es difícil separar la verdad de la fantasía,\\
la realidad de la imaginación.

A veces sospechas que la mitad de tus hazañas\\
no son más que reconstrucciones, tan lindas como tu esclava,\\
de algo tan simple y sucio y brutal,\\
conjuradas para cubrir la sangre y la mugre y el estiércol.

¿Pero a quién le importa ya? Os reunís cada año,\\
y cada año hay menos y menos invitados en casa de Menelao,\\
y coméis, y bebéis, y cantáis, para celebrar\\
el fingimiento de que vuestras vidas tuvieron algún sentido.

II

Cada vez menos. La mayoría ya no está.\\
Muchos de tus amigos no volvieron de Troya,\\
Aquiles, Áyax, Patroclo, Diomedes,\\
la lista es larga. Otros murieron en el viaje de vuelta.

Mucho se ha escrito sobre los dioses\\
oponiéndose al regreso de los argivos,\\
y poco sobre las ratas que, sospechas,\\
estaban detrás de las enfermedades que mataron a tantos.

Ah, pero a la gente le encanta cantar a los dioses,\\
ruines y caprichosos como son,\\
y celebrar sus enredos en los asuntos humanos,\\
su lujuria y su crueldad dementes.

Puestos a elegir, tú te quedarías con las ratas,\\
que, al menos, iban a lo suyo,\\
y aunque se comían tu comida y se bebían tu agua,\\
no se divertían con tu ruina, como los olímpicos.

III

Hace ya mucho tiempo\\
que dejaste de creer en esa gente\\
supuestamente divina, que bebía néctar y daba fiestas\\
en lo alto del monte Olimpo.

Lo más probable es que ninguno de tus compañeros,\\
los pocos viejos supervivientes aún en pie,\\
creyera en ellos. Pero todos fingís\\
una devoción conveniente para mantener el statu quo.

Porque, si los dioses no son reales,\\
¿por qué habría la gente de tomarse en serio a los reyes?\\
¿Por qué sufrir a tiranos como vosotros,\\
si están tan vacíos de sentido como las deidades?

Así que seguís con las quemas de muslos,\\
y las libaciones, y las plegarias, diligentemente conducidas\\
por esa otra especie de ratas, los sacerdotes,\\
siempre listos para comerse tu comida y beberse tu vino.

A cambio, por supuesto, de profecías\\
que resultan coincidir con tus propios planes,\\
y señales que confirman tus antojos,\\
y una bendición general de tus caprichos.

Así que, es verdad, los dioses no existen,\\
ni son necesarios, después de todo:\\
los hombres pueden encargarse de abusar,\\
y esclavizar, y matarse unos a otros.

IV

Los que todavía se reúnen cada año\\
fingen ser hermanos de sangre,\\
y es cierto, hay calor entre vosotros,\\
algo parecido a la amistad.

Ah, pero con qué avidez\\
calibra cada invitado el estado del otro,\\
mide la piel arruinada,\\
el pelo menguante, la espalda rota.

Con qué precisión evaluáis\\
los ojos apagados, la mano temblorosa,\\
los párpados caídos, el temblor en la voz,\\
todas las señales que gritan vejez.

Este ha engordado demasiado, este está demasiado flaco,\\
Idomeneo se pasa el día durmiendo, Filoctetes siempre está borracho,\\
Néstor apenas puede hablar, la mayoría\\
está en peor forma que tú.

Y te preguntas si el único sentido de esas reuniones\\
es reinventar el pasado,\\
espantar la soledad que avanza,\\
y contar los cuerpos supervivientes.

V

Menelao se conserva bastante bien,\\
puede argumentarse que él y tú\\
sois los menos dañados por el tiempo,\\
sobre todo porque ambos seguís lúcidos.

Y él es un verdadero hermano,\\
hay algo de verdad entre las muchas mentiras,\\
fuisteis hermanos de armas,\\
y más de una vez os salvasteis la vida el uno al otro.

Así que el vínculo es real,\\
tan real, de hecho, que solo Telémaco\\
es más querido para tu corazón,\\
sin contar, por supuesto, a Circe.

Pero tú no le mencionarás la hechicera\\
a tu amigo, como él no mencionará a Helena,\\
los dos habéis aprendido\\
a respetar el corazón roto del otro.

No es que seáis adolescentes enamorados,\\
las penas de los jóvenes se arreglan tan fácilmente\\
como sus cuerpos fuertes. Pero el amor perdido de un viejo\\
es el último dolor que llevará a la tumba.

VI

Y así, Menelao no habla de Helena,\\
hace ya tantos años que se fue,\\
de hecho desapareció poco después\\
de la visita de Telémaco a Esparta.

Tu amigo presumió tanto ante tu hijo,\\
se aseguró de que viera a la legendaria Helena en todo su esplendor,\\
para que se lo contara a su padre Odiseo,\\
por si acaso había sobrevivido.

Y en efecto, Telémaco dijo\\
que nunca había visto a nadie más hermoso,\\
ni más distante. Mirando su fría majestad,\\
se convenció de su naturaleza divina.

Así que cuando Helena se desvaneció, el relato oficial\\
habló de una santa ascensión al Elíseo,\\
no se había fugado otra vez, no,\\
sino que los dioses la habían reclamado.

VII

Nadie cuestionó la historia,\\
como nadie cuestiona tu aventura con el cíclope,\\
los cuentos se inventaron precisamente para esto,\\
para encontrar sentido donde solo hay ruido.

Ruido y furia, como las hordas que condujiste,\\
rabia, y odio y lujuria, bombeando en vuestros corazones,\\
bronce, y sangre y fuego,\\
para ser disfrazados de batallas épicas.

Así que fuisteis todos a Troya,\\
a vengar el rapto de la esposa de Menelao,\\
y luchasteis diez años, y quemasteis la ciudad,\\
y devolvisteis el trofeo al palacio.

Pero la segunda vez que ella desapareció,\\
no había ciudad que saquear, ni unidad entre los argivos,\\
y demasiados viejos al mando,\\
así que invocar a los dioses resultó apropiado.

VIII

Lo asombroso\\
es que Menelao la amaba.\\
La amó mucho más después de Troya\\
de lo que la había amado nunca antes.

Después de la boda,\\
Helena era solo un trofeo, y que el príncipe,\\
tan joven y furioso, se llevara a la cama a otras mujeres\\
era solo buen deporte, legal, para un hombre.

El supuesto secuestro lo enfureció,\\
quizás porque sabía, como todo el mundo,\\
que nadie había obligado a Helena a saltar\\
a la nave de Paris, ni a la cama de Paris.

Cuentan los relatos que en Troya\\
Helena desnudó el pecho ante su esposo,\\
lo invitó a golpear, y Menelao\\
dejó caer la espada y lloró.

Y por una vez estás convencido\\
de que esta historia sí es cierta:\\
Menelao había necesitado un rechazo y una guerra\\
para entender que estaba enamorado.

IX

Pero entonces cometió el error clásico:\\
asumió que su amor recién descubierto,\\
su deseo ferviente, su admiración por ella,\\
le ganarían el corazón.

Se equivocaba. Helena, por supuesto,\\
se comportaría como la noble dama que era,\\
y se ocuparía de la casa y del lecho,\\
guardándose sus cartas.

Pero fue su esposo quien mató a Deífobo,\\
y tú la viste llorar con el alma rota sobre su cuerpo,\\
qué ironía, pensaste. Ni el valiente Menelao,\\
ni el cobarde Paris, sino aquel otro,

príncipe discreto y valeroso,\\
casi tan bravo y honesto como Héctor.\\
A él, y solo a él,\\
le dio Helena su corazón.

X

Y por él, y solo por él,\\
tramó su venganza.\\
No eligió el cuchillo, como su hermana,\\
ni el veneno que tan bien conocía.

No: ella quería herir más hondo,\\
así que trabajó duro hasta que Menelao\\
solo podía respirar el aire que ella había respirado,\\
y entonces lo apuñaló por la espalda.

No hizo falta bronce, ni siquiera las nuevas dagas de hierro.\\
Ella, simplemente, se fue;\\
dicen algunos que un troyano llamado Eneas\\
organizó la expedición de rescate.

XI

Menelao no habla de ella,\\
como tú no hablas de Circe,\\
cada año, los dos os reunís\\
en Esparta y bebéis vino espeso,

y os emborracháis, y lloráis por el amigo perdido,\\
y por el amor perdido\\
que nunca merecisteis.
\end{verse}

\end{document}
